% ============================================================================
\chapter{Eisenstein 级数}
\label{ch:eisenstein}
% ============================================================================
% 对应 Diamond \& Shurman Chapter 4（在本书中为第 2 章）

在第一章末尾，我们见到了两类模形式的具体例子：
Eisenstein 级数 $E_4, E_6$ 和判别形式 $\Delta$。
本章的任务是深入研究 Eisenstein 级数——不仅要知道它"长什么样"，
还要\textbf{证明}它满足模变换律，并理解其系数的算术含义。

核心工具是\textbf{Poisson 求和公式}。
它把对格点的求和转化为对 Fourier 系数的求和，
是推导 $q$-展开的关键。Poisson 求和的思想非常普遍，
也是 Riemann $\zeta$ 函数、Dirichlet $L$-函数满足函数方程的根本原因。

本章结构：
\begin{itemize}
  \item \textbf{§\ref{sec:bernoulli}} Bernoulli 数与 $\zeta$ 函数特殊值——
    为 $q$-展开的常数项作准备。
  \item \textbf{§\ref{sec:lattice-sum-eisenstein}}
    从格求和出发定义 Eisenstein 级数，验证模变换律。
  \item \textbf{§\ref{sec:poisson-summation}}
    Poisson 求和公式及其证明。
  \item \textbf{§\ref{sec:eisenstein-q-expansion}}
    用 Poisson 求和推导 $E_k$ 的 $q$-展开，计算系数。
  \item \textbf{§\ref{sec:eisenstein-congruence}}
    同余子群上的 Eisenstein 级数与 Dirichlet 特征。
  \item \textbf{§\ref{sec:eisenstein-exercises}} 习题。
\end{itemize}


% ============================================================================
\section{Bernoulli 数与 $\zeta$ 函数}
\label{sec:bernoulli}
% ============================================================================

Eisenstein 级数 $E_k$ 的 $q$-展开常数项（即 $a_0$）
是 Riemann $\zeta$ 函数在偶数负整数处的特殊值，
可以用 Bernoulli 数表达。我们先把这些数字算清楚。

\textbf{Bernoulli 数是什么？}
把函数 $\frac{t}{e^t - 1}$ 展开成 $t$ 的幂级数：
\[
  \frac{t}{e^t - 1} = \sum_{m=0}^{\infty} \frac{B_m}{m!} t^m.
\]
这样定义的系数 $B_m$ 就叫\textbf{Bernoulli 数}\index{Bernoulli 数 (Bernoulli numbers)}。

\begin{example}[计算前几个 Bernoulli 数]
\label{ex:bernoulli-numbers}
用 $\frac{t}{e^t - 1} \cdot (e^t - 1) = t$ 可以逐步算出 $B_m$。
展开 $e^t - 1 = t + t^2/2! + t^3/3! + \cdots$，则
\[
  \left(\sum_{m=0}^{\infty} \frac{B_m}{m!} t^m\right)
  \left(\sum_{n=1}^{\infty} \frac{t^n}{n!}\right) = t.
\]
比较两边 $t^n$ 的系数（$n \geq 2$）：
\[
  \sum_{j=0}^{n-1} \frac{B_j}{j!} \cdot \frac{1}{(n-j)!} = 0,
  \qquad \text{即} \quad
  \sum_{j=0}^{n-1} \binom{n}{j} B_j = 0.
\]
从 $B_0 = 1$ 出发逐步计算：
\begin{align*}
  B_0 &= 1, \\
  B_1 &= -\frac{1}{2}, \\
  B_2 &= \frac{1}{6}, \quad B_4 = -\frac{1}{30}, \quad
  B_6 = \frac{1}{42}, \quad B_8 = -\frac{1}{30}, \quad
  B_{10} = \frac{5}{66}, \quad B_{12} = -\frac{691}{2730}.
\end{align*}
奇数阶（$m \geq 3$）：$B_m = 0$，因为
$\frac{t}{e^t-1} + \frac{t}{2} = \frac{t}{2} \cdot \coth\frac{t}{2}$
是偶函数。
\end{example}

注意 $B_{12} = -691/2730$——分子 $691$ 是个素数，
将在 Ramanujan $\tau$ 函数的同余式 $\tau(n) \equiv \sigma_{11}(n) \pmod{691}$
中再次出现。这不是巧合，背后有深刻的数论原因（Serre-Swinnerton-Dyer 理论）。

\begin{proposition}[Bernoulli 数与 $\zeta$ 函数]
\label{prop:bernoulli-zeta}
\index{Riemann zeta 函数 (Riemann zeta function)}\index{zeta@$\zeta(s)$}
对偶数 $k \geq 2$，
\[
  \zeta(k) = \sum_{n=1}^{\infty} \frac{1}{n^k}
  = -\frac{(2\pi i)^k B_k}{2 \cdot k!}.
\]
特别地，
\[
  \zeta(2) = \frac{\pi^2}{6}, \quad
  \zeta(4) = \frac{\pi^4}{90}, \quad
  \zeta(6) = \frac{\pi^6}{945}, \quad
  \zeta(8) = \frac{\pi^8}{9450}.
\]
\end{proposition}

\begin{proof}
利用 $\pi \cot(\pi z)$ 的部分分式展开。
由复分析，$\pi \cot(\pi z)$ 在所有整数处有单极点，留数均为 $1$，故
\[
  \pi \cot(\pi z) = \frac{1}{z} + \sum_{n=1}^{\infty}
  \left(\frac{1}{z-n} + \frac{1}{z+n}\right)
  = \frac{1}{z} + \sum_{n=1}^{\infty} \frac{2z}{z^2 - n^2}.
\]
另一方面，令 $q = e^{2\pi i z}$（$\im z > 0$），则
\[
  \pi\cot(\pi z)
  = \pi i \cdot \frac{e^{i\pi z} + e^{-i\pi z}}{e^{i\pi z} - e^{-i\pi z}}
  = \pi i \cdot \frac{q + 1}{q - 1}
  = -\pi i \left(1 + \frac{2q}{1-q}\right)
  = -\pi i - 2\pi i \sum_{n=1}^{\infty} q^n.
\]
比较两个展开在 $z = 0$ 附近的 $z^{2k-1}$（$k \geq 1$）系数，注意
\[
  \frac{2z}{z^2 - n^2} = -\frac{2}{n^2} \cdot \frac{1}{1 - (z/n)^2}
  = -\frac{2}{n^2} \sum_{j=0}^{\infty} \left(\frac{z}{n}\right)^{2j},
\]
故 $z^{2k-1}$ 系数贡献为 $-2\sum_{n=1}^\infty n^{-2k}$；
另一方面，$-2\pi i \sum_{n=1}^{\infty} e^{2\pi inz}$ 在 $z = 0$ 附近展开，
$z^{2k-1}$ 系数为 $-2\pi i \sum_{n=1}^\infty \frac{(2\pi i n)^{2k-1}}{(2k-1)!}$。
整理后即得 $\zeta(2k) = -\frac{(2\pi i)^{2k} B_{2k}}{2(2k)!}$。
\end{proof}

\textbf{系数的化简。}
由\cref{prop:bernoulli-zeta}，
\[
  \frac{2}{2\zeta(k)} = \frac{2}{-\frac{(2\pi i)^k B_k}{k!}}
  = \frac{-2k!}{(2\pi i)^k B_k}.
\]
这在后面计算 $E_k$ 的 $q$-展开系数时会用到。


% ============================================================================
\section{格求和与 Eisenstein 级数的定义}
\label{sec:lattice-sum-eisenstein}
% ============================================================================

在\cref{sec:modular-forms-def}中我们预告了 Eisenstein 级数。
现在我们从头构造它，并直接验证模变换律。

\textbf{动机：用格求和构造自动满足变换律的函数。}
模形式的变换律是 $f(\gamma(\tau)) = (c\tau + d)^k f(\tau)$。
如果我们对某个函数 $h$ 做"平均"：
\[
  f(\tau) = \sum_{\gamma \in \SL_2(\Z)} h(\gamma(\tau)) \cdot j(\gamma, \tau)^{-k},
  \qquad j(\gamma, \tau) = c\tau + d,
\]
那么 $f(\gamma_0(\tau)) = j(\gamma_0, \tau)^k f(\tau)$ 自动成立——
因为 $\gamma \mapsto \gamma\gamma_0$ 只是对求和指标做了置换。
Eisenstein 级数就是取最简单的 $h(\tau) = 1$ 时得到的结果。

\begin{definition}[Eisenstein 级数 $G_k$]
\label{def:eisenstein-Gk}
\index{Eisenstein 级数 (Eisenstein series)}\index{Gk@$G_k(\tau)$}
对偶整数 $k \geq 4$，定义
\[
  G_k(\tau) = \sum_{\substack{(c,d) \in \Z^2 \\ (c,d) \neq (0,0)}}
  \frac{1}{(c\tau + d)^k}.
\]
\end{definition}

\begin{example}[前几项展开]
\label{ex:Gk-first-terms}
把 $G_k(\tau)$ 的求和按 $c$ 的值分组：
\begin{itemize}
  \item $c = 0$：贡献 $\sum_{d \neq 0} d^{-k} = 2\zeta(k)$。
  \item $c \neq 0$：对固定 $c$，$d$ 跑遍所有整数，
    贡献 $\sum_{d \in \Z} (c\tau + d)^{-k}$。
\end{itemize}
故
\[
  G_k(\tau) = 2\zeta(k) + 2 \sum_{c=1}^{\infty} \sum_{d \in \Z} \frac{1}{(c\tau + d)^k}.
\]
\end{example}

\begin{proposition}[绝对收敛性]
\label{prop:Gk-convergence}
对 $k \geq 3$，$G_k(\tau)$ 在 $\HH$ 上绝对收敛，
且在每个紧子集上一致收敛。
\end{proposition}

\begin{proof}
设 $\tau = x + iy$，$y > 0$。对 $(c, d) \neq (0, 0)$，
\[
  |c\tau + d|^2 = (cx + d)^2 + c^2 y^2 \geq \min(1, y^2)(c^2 + d^2).
\]
（若 $c = 0$：$|d|^2 = d^2 \geq c^2+d^2$；若 $c \neq 0$：
$|c\tau+d|^2 \geq c^2y^2 \geq y^2(c^2+d^2)/2$ 当 $|c| \leq |d|$，
直接估计另一情形。）

故
\[
  |G_k(\tau)| \leq C(y) \sum_{(c,d) \neq 0} (c^2 + d^2)^{-k/2}.
\]
对二维格求和，当 $k > 2$ 时级数收敛（在极坐标下估计）。
在紧子集 $\{y \geq \delta\}$ 上 $C(y)$ 有上界，故一致收敛。
\end{proof}

\begin{theorem}[$G_k$ 是模形式]
\label{thm:Gk-modular}
对偶整数 $k \geq 4$，$G_k(\tau)$ 是 $\SL_2(\Z)$ 的权 $k$ 模形式，
且 $G_k(\infty) = 2\zeta(k) \neq 0$。
\end{theorem}

\begin{proof}
\textbf{全纯性}：由\cref{prop:Gk-convergence}，一致收敛的全纯函数之极限是全纯的。

\textbf{变换律}：设 $\gamma_0 = \begin{pmatrix} a & b \\ c_0 & d_0 \end{pmatrix}
\in \SL_2(\Z)$，$\gamma_0(\tau) = (a\tau + b)/(c_0\tau + d_0)$。
\begin{align}
  G_k(\gamma_0(\tau))
  &= \sum_{(c,d) \neq 0} \frac{1}{\!\left(c \cdot \frac{a\tau+b}{c_0\tau+d_0} + d\right)^{\!k}} \\
  &= (c_0\tau + d_0)^k
     \sum_{(c,d) \neq 0} \frac{1}{((ca + dc_0)\tau + (cb + dd_0))^k}.
\end{align}
映射 $(c, d) \mapsto (ca+dc_0,\; cb+dd_0)$ 是 $\Z^2 \setminus \{0\}$ 上的双射
（由 $\det\gamma_0 = 1$），故
\[
  G_k(\gamma_0(\tau)) = (c_0\tau + d_0)^k G_k(\tau).
\]

\textbf{在尖点处}：$\im(\tau) \to +\infty$ 时非常数项趋于 $0$
（将在\cref{sec:eisenstein-q-expansion}中用 Poisson 求和证明），
故 $G_k(\tau) \to 2\zeta(k)$，在尖点处全纯。
\end{proof}

\begin{definition}[规范化 Eisenstein 级数 $E_k$]
\label{def:normalized-eisenstein}
\index{Ek@$E_k(\tau)$}
\[
  E_k(\tau) = \frac{G_k(\tau)}{2\zeta(k)},
\]
使得 $E_k(\infty) = 1$。
\end{definition}

\begin{example}[$E_4, E_6$ 的展开（预告）]
\label{ex:E4-E6-expansion}
用 Poisson 求和（\cref{sec:eisenstein-q-expansion}）可以证明：
\begin{align*}
  E_4(\tau) &= 1 + 240\sum_{n=1}^{\infty} \sigma_3(n) q^n
  = 1 + 240q + 2160q^2 + 6720q^3 + \cdots \\
  E_6(\tau) &= 1 - 504\sum_{n=1}^{\infty} \sigma_5(n) q^n
  = 1 - 504q - 16632q^2 - 122976q^3 - \cdots
\end{align*}
其中 $\sigma_{k-1}(n) = \sum_{d \mid n} d^{k-1}$。
系数 $240 = -2k/B_k|_{k=4} = -8/(-1/30) = 240$ ($\checkmark$)，
$-504 = -2 \cdot 6 / B_6 = -12 / (1/42) = -504$ ($\checkmark$)。
\end{example}

% figures/ch04-eisenstein-series-plot.tex
% Eisenstein 级数 E_4 的 q-展开前几项示意（数轴上的系数增长）
\begin{figure}[ht]
\centering
\begin{tikzpicture}[>=Stealth, font=\small, scale=1.0]

% 坐标轴
\draw[->] (-0.3, 0) -- (7.5, 0) node[right] {$n$};
\draw[->] (0, -0.3) -- (0, 4.2) node[above] {$a_n(E_4)$（规范化后）};

% 刻度
\foreach \x in {1,2,3,4,5,6,7} {
  \draw (\x, 0.05) -- (\x, -0.05);
  \node[below] at (\x, -0.08) {\x};
}

% 系数值（规范化：a_0=1, a_n = 240 σ_3(n) / 240 = σ_3(n)，再缩放显示）
% a_0=1, a_1=240, a_2=2160, a_3=6720, a_4=17520, a_5=20160, a_6=60480
% 缩放：除以 240 * 5 = 1200 再乘以 3.5 的高度
% σ_3(1)=1, σ_3(2)=9, σ_3(3)=28, σ_3(4)=73, σ_3(5)=126, σ_3(6)=252

% 用柱状图
\foreach \x/\h/\lab in {
  1/0.35/{$1$},
  2/3.15/{$9$},
  3/3.99/{$28$（峰值截断）},
  4/3.8/{$73$（截断）},
  5/3.9/{$126$（截断）},
  6/4.0/{$252$（截断）},
  7/3.85/{$344$（截断）}
}{
  \fill[blue!40!white, draw=blue!60!black]
    (\x-0.3, 0) rectangle (\x+0.3, \h);
}

% a_0 = 1（常数项）
\fill[orange!60!white, draw=orange!80!black]
  (-0.3, 0) rectangle (0.3, 0.35);
\node[below, orange!70!black] at (0, -0.08) {$0$};

% 标注
\node[blue!70!black, above right, font=\footnotesize] at (1, 0.35) {$\sigma_3(1)=1$};
\node[blue!70!black, above, font=\footnotesize] at (2, 3.15) {$\sigma_3(2)=9$};
\node[blue!70!black, above right, font=\footnotesize] at (6, 4.0) {$\sigma_3(n)\nearrow$};

% 公式提示
\node[font=\footnotesize, align=left, text=gray!70!black] at (5.5, 1.5)
  {$E_4(\tau) = 1 + 240\!\sum_{n=1}^{\infty}\!\sigma_3(n)\,q^n$};

% 虚线截断线
\draw[dashed, gray!60] (0, 4.0) -- (7.5, 4.0);
\node[right, font=\footnotesize, gray!60] at (7.5, 4.0) {（显示截断）};

\end{tikzpicture}
\caption{$E_4$ 的 $q$-展开系数 $a_n = 240\sigma_3(n)$（$n \geq 1$）随 $n$ 增长的示意。
因子 $\sigma_3(n) = \sum_{d \mid n} d^3$ 是因子的立方和，增长迅速。
系数的精确公式由 Poisson 求和公式导出。}
\label{fig:eisenstein-coefficients}
\end{figure}



% ============================================================================
\section{Poisson 求和公式}
\label{sec:poisson-summation}
% ============================================================================

Poisson 求和公式是推导 $E_k$ 的 $q$-展开的核心工具，
也是联系格求和与 Fourier 分析的桥梁。
它的思想很简单：对一个函数"按整数周期化"，
有两种方式——直接加格点，或者在频域加频率——结果相同。

\begin{definition}[Fourier 变换]
\label{def:fourier-transform}
\index{Fourier 变换 (Fourier transform)}
设 $f\colon \R \to \C$ 是 Schwartz 函数（光滑且各阶导数迅速衰减）。
其 \textbf{Fourier 变换}定义为
\[
  \hat{f}(\xi) = \int_{-\infty}^{+\infty} f(x)\, e^{-2\pi i \xi x}\, dx.
\]
\end{definition}

\begin{example}[高斯函数的 Fourier 变换]
\label{ex:gaussian-fourier}
\index{高斯函数 (Gaussian function)}
取 $f(x) = e^{-\pi x^2}$（高斯函数，满足 $\hat{f} = f$，
这是 Schwartz 函数中最特别的一个）。直接计算：
\[
  \hat{f}(\xi)
  = \int_{-\infty}^{+\infty} e^{-\pi x^2} e^{-2\pi i \xi x}\, dx.
\]
配方：$-\pi x^2 - 2\pi i \xi x = -\pi(x + i\xi)^2 - \pi\xi^2$，故
\[
  \hat{f}(\xi) = e^{-\pi \xi^2} \int_{-\infty}^{+\infty} e^{-\pi(x+i\xi)^2}\, dx.
\]
通过围道积分（围绕一个矩形区域），积分值等于
$\int_{-\infty}^{+\infty} e^{-\pi x^2}\, dx = 1$（高斯积分）。
所以 $\hat{f}(\xi) = e^{-\pi\xi^2} = f(\xi)$：
\textbf{高斯函数是自己的 Fourier 变换！}

稍作推广：对 $a > 0$，取 $f_a(x) = e^{-\pi a x^2}$，则
$\hat{f}_a(\xi) = a^{-1/2} e^{-\pi \xi^2/a}$。
\end{example}

\begin{theorem}[Poisson 求和公式]
\label{thm:poisson-summation}
\index{Poisson 求和公式 (Poisson summation formula)}
设 $f\colon \R \to \C$ 是 Schwartz 函数（或更一般地，$f$ 和 $\hat{f}$ 均绝对可积且连续）。则
\[
  \sum_{n \in \Z} f(n) = \sum_{n \in \Z} \hat{f}(n).
\]
\end{theorem}

% figures/ch04-poisson-summation.tex
% Poisson 求和公式示意：整数格点求和 ↔ Fourier 频率求和
\begin{figure}[ht]
\centering
\begin{tikzpicture}[>=Stealth, font=\small, scale=1.0]

% === 左侧：f 在整数点的求和 ===
\begin{scope}[xshift=-4cm]
  \draw[->] (-3.2, 0) -- (3.2, 0) node[right] {$x$};
  \draw[->] (0, -0.2) -- (0, 2.5) node[above] {$f(x)$};

  % 连续函数曲线（高斯型）
  \draw[thick, blue!70!black, domain=-3:3, samples=80]
    plot (\x, {2.0*exp(-0.7*\x*\x)});

  % 整数点上的线段（柱高度硬编码）
  % n=-3: h=2e^{-6.3}≈0.04, n=-2: h=2e^{-2.8}≈0.12, n=-1: h=2e^{-0.7}≈0.99
  % n=0: h=2, n=1: h≈0.99, n=2: h≈0.12, n=3: h≈0.04
  \foreach \n/\h in {-3/0.04, -2/0.12, -1/0.99, 0/2.0, 1/0.99, 2/0.12, 3/0.04}{
    \draw[red!80!black, line width=2pt] (\n, 0) -- (\n, \h);
    \fill[red!80!black] (\n, \h) circle (2.5pt);
  }

  \node[below, font=\footnotesize] at (0, -0.5)
    {$\displaystyle\sum_{n \in \Z} f(n)$};
  \node[above right, blue!70!black, font=\footnotesize] at (1.5, 1.6)
    {$f(x)$};
\end{scope}

% === 中间等号 ===
\node[font=\Large] at (0, 1) {$=$};
\node[font=\footnotesize, text=gray!70!black, align=center] at (0, 0.2)
  {Poisson\\求和};

% === 右侧：hat{f} 在整数点的求和 ===
\begin{scope}[xshift=4cm]
  \draw[->] (-3.2, 0) -- (3.2, 0) node[right] {$\xi$};
  \draw[->] (0, -0.2) -- (0, 2.5) node[above] {$\hat{f}(\xi)$};

  % Fourier变换后的高斯（宽度不同）
  \draw[thick, green!50!black, domain=-3:3, samples=80]
    plot (\x, {1.7*exp(-1.4*\x*\x)});

  % 整数点: n=0: 1.7, n=±1: 1.7e^{-1.4}≈0.42, n=±2: 1.7e^{-5.6}≈0.007
  \foreach \n/\h in {-2/0.01, -1/0.42, 0/1.7, 1/0.42, 2/0.01}{
    \draw[orange!80!black, line width=2pt] (\n, 0) -- (\n, \h);
    \fill[orange!80!black] (\n, \h) circle (2.5pt);
  }

  \node[below, font=\footnotesize] at (0, -0.5)
    {$\displaystyle\sum_{n \in \Z} \hat{f}(n)$};
  \node[above right, green!50!black, font=\footnotesize] at (1.2, 1.2)
    {$\hat{f}(\xi)$};
\end{scope}

% === 底部公式 ===
\node[font=\footnotesize, align=center] at (0, -1.4)
  {$\hat{f}(\xi) = \int_{-\infty}^{+\infty} f(x)\,e^{-2\pi i\xi x}\,dx$};

\end{tikzpicture}
\caption{Poisson 求和公式的直觉。对光滑衰减函数 $f$，
在整数格点上的求和 $\sum_{n \in \Z} f(n)$（左，红色竖线）
等于其 Fourier 变换 $\hat{f}$ 在整数点上的求和 $\sum_{n \in \Z} \hat{f}(n)$（右，橙色竖线）。
关键步骤：把 $\sum_n f(n+x)$ 展开为 Fourier 级数，再令 $x=0$。}
\label{fig:poisson-summation}
\end{figure}


\begin{proof}
定义周期函数
\[
  F(x) = \sum_{n \in \Z} f(x + n).
\]
（由 $f$ 的 Schwartz 性质，对每个 $x$ 这个级数绝对收敛，
且 $F$ 是周期为 $1$ 的光滑函数。）

$F$ 有 Fourier 级数展开：
\[
  F(x) = \sum_{m \in \Z} c_m e^{2\pi i mx},
  \qquad c_m = \int_0^1 F(x) e^{-2\pi imx}\, dx.
\]
计算 $c_m$：
\begin{align}
  c_m
  &= \int_0^1 \left(\sum_{n \in \Z} f(x+n)\right) e^{-2\pi imx}\, dx \\
  &= \sum_{n \in \Z} \int_0^1 f(x+n) e^{-2\pi imx}\, dx \\
  &= \sum_{n \in \Z} \int_n^{n+1} f(t) e^{-2\pi im(t-n)}\, dt
    \qquad (t = x + n,\; e^{2\pi imn} = 1) \\
  &= \int_{-\infty}^{+\infty} f(t) e^{-2\pi imt}\, dt = \hat{f}(m).
\end{align}
（交换求和与积分由绝对收敛保证；第三步用了 $e^{2\pi imn} = 1$。）

故 $F(x) = \sum_{m \in \Z} \hat{f}(m) e^{2\pi imx}$。
令 $x = 0$：
\[
  F(0) = \sum_{n \in \Z} f(n) = \sum_{m \in \Z} \hat{f}(m). \qedhere
\]
\end{proof}

\textbf{推广：$t > 0$ 的参数化版本。}
设 $g_t(x) = f(tx)$，则 $\hat{g}_t(\xi) = t^{-1}\hat{f}(\xi/t)$。
应用 Poisson 求和：
\[
  \sum_{n \in \Z} f(nt) = \frac{1}{t} \sum_{n \in \Z} \hat{f}(n/t).
\]
这在推导 $\theta$ 函数的变换律时特别有用。

\begin{example}[高斯求和与 $\theta$ 函数]
\label{ex:theta-poisson}
\index{theta 函数 (theta function)}\index{θ@$\theta(t)$}
取 $f(x) = e^{-\pi x^2}$，$\hat{f}(\xi) = e^{-\pi\xi^2}$。
定义 $\theta(t) = \sum_{n \in \Z} e^{-\pi n^2 t}$（$t > 0$）。
用参数化 Poisson 求和（$f \to f(x)$，$t \to \sqrt{t}$）：
\[
  \theta(t) = \sum_{n \in \Z} e^{-\pi n^2 t}
  = \frac{1}{\sqrt{t}} \sum_{n \in \Z} e^{-\pi n^2/t}
  = \frac{1}{\sqrt{t}} \theta(1/t).
\]
这就是 $\theta$ 函数的\textbf{函数方程}，
也是 Riemann $\zeta$ 函数满足对称函数方程的根源：
\[
  \theta(t) = \frac{1}{\sqrt{t}} \theta(1/t)
  \quad \Longrightarrow \quad
  \xi(s) = \xi(1-s),
  \quad \xi(s) = \frac{1}{2}s(s-1)\pi^{-s/2}\Gamma(s/2)\zeta(s).
\]
\end{example}

\begin{example}[用 Poisson 求和计算 $\sum_{d \in \Z} (c\tau + d)^{-k}$]
\label{ex:poisson-for-eisenstein}
这是推导 $G_k$ 的 $q$-展开的直接预备。
固定 $c \in \Z_{>0}$，$\tau \in \HH$，计算
\[
  S_c(\tau) = \sum_{d \in \Z} \frac{1}{(c\tau + d)^k}.
\]
取 $f(x) = (c\tau + x)^{-k}$，$\hat{f}(\xi) = \int_{-\infty}^{+\infty} (c\tau + x)^{-k} e^{-2\pi i \xi x}\, dx$。

令 $z = c\tau + x$（$\im(c\tau) = c\, \im\tau > 0$），$dx = dz$：
\[
  \hat{f}(\xi)
  = e^{-2\pi i \xi(c\tau - c\tau)} \int_{\im(c\tau) + \R} z^{-k} e^{-2\pi i\xi z}\, dz
  = e^{2\pi i \xi c\tau} \int_{\R + i c\, \im\tau} z^{-k} e^{-2\pi i \xi z}\, dz.
\]

\textbf{分两种情形：}

\textbf{情形 1（$\xi \leq 0$）}：当 $\xi \leq 0$ 时，
$e^{-2\pi i\xi z}$ 在上半平面（$\im z > 0$）没有奇点，
围道可以推到上方无穷远而积分趋于零。
故 $\hat{f}(\xi) = 0$（$\xi \leq 0$）。

\textbf{情形 2（$\xi > 0$）}：$e^{-2\pi i\xi z}$ 在下半平面衰减，
围道推到下方，通过 $z = 0$ 处的极点。由留数定理：
\[
  \int z^{-k} e^{-2\pi i\xi z}\, dz = -2\pi i \cdot \Res_{z=0}\left(z^{-k} e^{-2\pi i\xi z}\right).
\]
在 $z = 0$ 附近展开 $e^{-2\pi i\xi z} = \sum_{j=0}^\infty \frac{(-2\pi i\xi)^j}{j!} z^j$，
则 $z^{-k} e^{-2\pi i\xi z}$ 在 $z^{-1}$ 处的系数（留数）为
$\frac{(-2\pi i\xi)^{k-1}}{(k-1)!}$，故
\[
  \hat{f}(\xi) = e^{2\pi i\xi c\tau} \cdot (-2\pi i) \cdot \frac{(-2\pi i\xi)^{k-1}}{(k-1)!}
  = \frac{(-2\pi i)^k \xi^{k-1}}{(k-1)!} e^{2\pi i\xi c\tau}.
\]

代入 Poisson 求和（注意 $\hat{f}(\xi) = 0$ 对 $\xi \leq 0$，对整数 $\xi = m \geq 1$）：
\[
  S_c(\tau) = \sum_{d \in \Z} f(d) = \sum_{m=1}^{\infty} \hat{f}(m)
  = \frac{(-2\pi i)^k}{(k-1)!} \sum_{m=1}^{\infty} m^{k-1} e^{2\pi i m c\tau}.
\]
\end{example}


% ============================================================================
\section{$q$-展开的推导}
\label{sec:eisenstein-q-expansion}
% ============================================================================

有了 Poisson 求和的准备，现在可以完整地推导 $G_k$ 的 $q$-展开。

% figures/ch04-q-expansion-geometry.tex
% q 展开的几何意义：τ → q = e^{2πiτ} 把竖带映到穿孔圆盘
\begin{figure}[ht]
\centering
\begin{tikzpicture}[>=Stealth, font=\small, scale=1.1]

% === 左侧：上半平面的竖带 ===
\begin{scope}[xshift=-4.2cm]
  % 阴影填充带
  \fill[blue!8]
    (-1, 0.3) rectangle (1, 3.5);
  % 边界
  \draw[thick, blue!60!black] (-1, 0.3) -- (-1, 3.5);
  \draw[thick, blue!60!black] (1, 0.3) -- (1, 3.5);
  \draw[dashed, gray] (-1.5, 0.3) -- (1.5, 0.3)
    node[right, font=\footnotesize] {$\im(\tau) = Y$};
  % 坐标轴
  \draw[->] (-1.6, 0) -- (1.6, 0) node[right] {$\Re(\tau)$};
  \draw[->] (0, 0) -- (0, 3.8) node[above] {$\im(\tau)$};
  % 刻度
  \node[below] at (-1, 0) {$-\tfrac{1}{2}$};
  \node[below] at (1, 0) {$\tfrac{1}{2}$};
  % 示例点
  \fill[red!70!black] (0.3, 1.5) circle (2pt);
  \node[right, red!70!black, font=\footnotesize] at (0.33, 1.5) {$\tau$};
  % 顶部箭头（∞方向）
  \draw[->, green!50!black, line width=1.2pt] (0, 3.2) -- (0, 3.7);
  \node[green!50!black, right, font=\footnotesize] at (0.05, 3.5) {$\infty$};
  % 标签
  \node[blue!60!black] at (0, 0.9) {$\HH$（竖带）};
\end{scope}

% === 中间箭头 ===
\draw[->, thick, purple!70!black, line width=1.5pt]
  (-1.8, 1.8) -- (-0.3, 1.8)
  node[above, midway, font=\footnotesize, purple!70!black]
  {$q = e^{2\pi i\tau}$};

% === 右侧：穿孔圆盘 ===
\begin{scope}[xshift=1.8cm]
  % 圆盘
  \fill[orange!10] (0,0) circle (2cm);
  \fill[white] (0,0) circle (0.25cm);  % 穿孔
  \draw[thick, orange!70!black] (0,0) circle (2cm);
  \draw[thick, dashed, orange!40!black] (0,0) circle (0.25cm);
  % r=|q|=e^{-2πY}的圆
  \draw[dashed, blue!50!black] (0,0) circle (1.2cm);
  \node[blue!50!black, font=\footnotesize] at (1.05, 0.65)
    {$|q|=e^{-2\pi Y}$};
  % 示例点 q
  \fill[red!70!black] (0.6, 0.9) circle (2pt);
  \node[right, red!70!black, font=\footnotesize] at (0.63, 0.9) {$q(\tau)$};
  % 穿孔标签
  \node[font=\footnotesize, gray] at (0, -0.45) {$q=0$};
  \draw[->, gray, shorten >=3pt] (0.0, -0.42) -- (0.0, -0.23);
  % 外圆标签
  \node[font=\footnotesize, orange!70!black] at (1.6, -1.6)
    {$|q|=1$};
  \node[font=\footnotesize] at (0, -2.5) {穿孔圆盘 $\mathbb{D}^*$};
\end{scope}

% === 对应关系注记 ===
\node[font=\footnotesize, align=left, text=gray!70!black] at (0, -1.5)
  {$\im(\tau)\to+\infty$ $\Leftrightarrow$ $q\to 0$\quad；\quad
   $\tau \equiv \tau+1$ $\Leftrightarrow$ $q$ 不变（周期性）};

\end{tikzpicture}
\caption{$q$-展开的几何意义。映射 $q = e^{2\pi i\tau}$ 把上半平面中的竖带
$\{|\Re(\tau)| \leq \frac{1}{2},\, \im(\tau) > Y\}$（左）
双全纯地映到穿孔圆盘 $\{0 < |q| < e^{-2\pi Y}\}$（右）。
$\tau \to \infty$ 对应 $q \to 0$，模形式在尖点全纯等价于在 $q=0$ 处可去奇点。}
\label{fig:q-expansion-geometry}
\end{figure}


\begin{theorem}[$G_k$ 的 $q$-展开]
\label{thm:Gk-q-expansion}
对偶整数 $k \geq 4$，令 $q = e^{2\pi i\tau}$（$\tau \in \HH$），则
\[
  G_k(\tau) = 2\zeta(k)
  + \frac{2(2\pi i)^k}{(k-1)!} \sum_{n=1}^{\infty} \sigma_{k-1}(n)\, q^n,
\]
其中 $\sigma_{k-1}(n) = \sum_{d \mid n} d^{k-1}$ 是 $n$ 的因子的 $(k-1)$ 次幂之和。
\end{theorem}

\begin{proof}
由\cref{ex:Gk-first-terms}，
\[
  G_k(\tau) = 2\zeta(k) + 2\sum_{c=1}^{\infty} S_c(\tau),
  \qquad S_c(\tau) = \sum_{d \in \Z} (c\tau + d)^{-k}.
\]
由\cref{ex:poisson-for-eisenstein}（Poisson 求和对 $S_c$），
\[
  S_c(\tau) = \frac{(-2\pi i)^k}{(k-1)!} \sum_{m=1}^{\infty} m^{k-1} e^{2\pi i mc\tau}
  = \frac{(2\pi i)^k}{(k-1)!} \sum_{m=1}^{\infty} m^{k-1} q^{mc}.
\]
（用了 $(-1)^k = 1$ 因为 $k$ 为偶数。）

代入：
\[
  G_k(\tau)
  = 2\zeta(k) + \frac{2(2\pi i)^k}{(k-1)!}
    \sum_{c=1}^{\infty} \sum_{m=1}^{\infty} m^{k-1} q^{mc}.
\]
交换求和顺序，令 $n = mc$（固定 $n$，$m$ 跑遍 $n$ 的正因子，$c = n/m$）：
\begin{align}
  \sum_{c=1}^{\infty} \sum_{m=1}^{\infty} m^{k-1} q^{mc}
  &= \sum_{n=1}^{\infty} \left(\sum_{m \mid n} m^{k-1}\right) q^n
  = \sum_{n=1}^{\infty} \sigma_{k-1}(n)\, q^n.
\end{align}
这就完成了证明。
\end{proof}

\textbf{从 $G_k$ 到 $E_k$。}
用\cref{prop:bernoulli-zeta}，$2\zeta(k) = -\frac{(2\pi i)^k B_k}{k!}$，故
\[
  E_k(\tau)
  = \frac{G_k(\tau)}{2\zeta(k)}
  = 1 - \frac{2k}{B_k} \sum_{n=1}^{\infty} \sigma_{k-1}(n)\, q^n.
\]
（验证：$\frac{2(2\pi i)^k/(k-1)!}{2\zeta(k)} = \frac{2(2\pi i)^k/(k-1)!}{-(2\pi i)^k B_k/k!}
= \frac{-2k!}{(k-1)! B_k} = \frac{-2k}{B_k}$。）

\begin{example}[$\sigma_{k-1}$ 的计算]
\label{ex:sigma-computation}
\leavevmode
\begin{itemize}
  \item $\sigma_3(1) = 1^3 = 1$，$\sigma_3(2) = 1^3 + 2^3 = 9$，
    $\sigma_3(3) = 1 + 27 = 28$，$\sigma_3(4) = 1 + 8 + 64 = 73$，
    $\sigma_3(6) = 1 + 8 + 27 + 216 = 252$。
  \item $\sigma_5(1) = 1$，$\sigma_5(2) = 1 + 32 = 33$，
    $\sigma_5(3) = 1 + 243 = 244$，$\sigma_5(4) = 1 + 32 + 1024 = 1057$。
\end{itemize}

验证 $E_4$ 的前几项：
$1 + 240(q + 9q^2 + 28q^3 + 73q^4 + \cdots)
= 1 + 240q + 2160q^2 + 6720q^3 + 17520q^4 + \cdots$ \quad $(\checkmark)$
\end{example}

\begin{proposition}[$E_4, E_6$ 与 $\Delta$]
\label{prop:E4-E6-Delta}
\index{判别形式 (discriminant form)}\index{Delta@$\Delta(\tau)$}
\index{Ramanujan tau 函数}\index{tau@$\tau(n)$}
定义
\[
  \Delta(\tau) = \frac{E_4(\tau)^3 - E_6(\tau)^2}{1728}.
\]
则 $\Delta$ 是权 $12$ 的\textbf{尖点形式}（$\Delta(\infty) = 0$），
且有乘积公式
\[
  \Delta(\tau) = q \prod_{n=1}^{\infty}(1-q^n)^{24}
  = \sum_{n=1}^{\infty} \tau(n)\, q^n,
\]
其中 $\tau(n)$ 是 \textbf{Ramanujan $\tau$-函数}，$\tau(1) = 1$，$\tau(2) = -24$，
$\tau(3) = 252$，$\tau(4) = -1472$，$\ldots$
\end{proposition}

\begin{proof}
\textbf{$\Delta$ 是尖点形式}：$E_4(\infty) = 1$，$E_6(\infty) = 1$，故
$\Delta(\infty) = (1 - 1)/1728 = 0$，满足尖点条件。
权 $12 = k$ 由 $E_4^3$（权 $4 \times 3 = 12$）和 $E_6^2$（权 $6 \times 2 = 12$）的差给出。

\textbf{$\Delta$ 在 $\HH$ 上无零点（乘积公式）}：这比较深刻，
用到模形式空间的维数理论（见\cref{ch:dimension}）：
$\Sk_{12}(\SL_2(\Z))$ 恰好是一维的，
且 $\eta(\tau)^{24}$ 是其中一个非零元（Dedekind $\eta$ 函数
$\eta(\tau) = q^{1/24}\prod_{n=1}^\infty(1-q^n)$ 的 $24$ 次幂），
由比较常数项得 $\Delta(\tau) = q\prod_{n=1}^\infty(1-q^n)^{24}$。
\end{proof}

\textbf{$\Delta$ 与 Ramanujan 猜想。}
Ramanujan 在 1916 年观察到 $\tau(n)$ 的三个猜想：
\begin{enumerate}
  \item \textbf{乘性}：$\tau(mn) = \tau(m)\tau(n)$ 当 $\gcd(m,n) = 1$；
  \item \textbf{Hecke 关系}：$\tau(p^{r+1}) = \tau(p)\tau(p^r) - p^{11}\tau(p^{r-1})$；
  \item \textbf{估计}：$|\tau(p)| \leq 2p^{11/2}$。
\end{enumerate}
前两条由 Mordell（1917）证明，第三条（Ramanujan 猜想）由 Deligne 在 1974 年
利用 Weil 猜想的证明给出。乘性与 Hecke 算子有关（第三章），
估计与代数几何有关，这两个方向最终都汇入费马大定理的证明链。

\begin{example}[验证 $\tau$ 的乘性]
\label{ex:tau-multiplicative}
验证 $\tau(2)\tau(3) = \tau(6)$：
$(-24)(252) = -6048$，而 $\tau(6) = -6048$（$\checkmark$，需要计算 $q^6$ 系数）。

详细：$\Delta = q - 24q^2 + 252q^3 - 1472q^4 + 4830q^5 - 6048q^6 + \cdots$
\end{example}
